\section{Related Work}~\label{related}
In this section, we review related work along three directions. We first discuss knowledge-based robotic systems for structured environment representation. We then review planning under uncertainty, including POMDP-based and belief space methods. Finally, we examine human corrective systems that incorporate feedback for decision making.



\begin{table*}[t]
\caption{Qualitative comparison of related works.}
\label{tab:comparison}
\centering
\small
\setlength{\tabcolsep}{3pt}
\begin{tabularx}{\textwidth}{p{0.25\textwidth}
>{\centering\arraybackslash}X
>{\centering\arraybackslash}X
>{\centering\arraybackslash}X
>{\centering\arraybackslash}X
>{\centering\arraybackslash}X}
\toprule
Related Works
& \makecell{Interpretable\\Knowledge}
& \makecell{Symbolic\\Transition}
& \makecell{Uncertainty-Aware\\Planning}
& \makecell{Human\\Feedback}
& \makecell{System-Initiated\\Query} \\
\midrule

Logic-based knowledge inference

~\cite{tenorth2012knowledge,liu2020graph,
kunze2014using,hasegawa2023inferring,
mota2021integrated,zhang2024icorpp,daruna2019robocse,
mitrevsk2021ontology,olivares2019review,manzoor2021ontology}
& $\bigcirc$
& $\times$
& △
& $\times$ 
& $\times$ \\

Classical planning

~\cite{fikes1971strips,aeronautiques1998pddl,
helmert2009concise,hoffmann2001ff,
helmert2006fast,correa2024lifted}
& $\bigcirc$
& $\bigcirc$
& $\times$
& $\times$ 
& $\times$ \\

LLM-based knowledge inference

~\cite{ahn2022can,huang2022inner,singh2023progprompt,liang2023code,ding2023task,gu2024conceptgraphs,leanza2025conceptbot}
& $\bigcirc$
& △
& $\times$
& $\times$
& $\times$ \\

LLM-based planning

~\cite{singh2025adaptbot,zhouenhancing,xu2024p,
wang2025instructrag,kumar2026open,
huang2023voxposer,dai2024optimal,
liu2025delta,paulius2025bootstrapping,
quartey2025verifiably}
& $\bigcirc$
& △
& △
& △
& $\times$ \\

POMDP-based planning

~\cite{somani2013despot,bai2015intention,cramer2021probabilistic,bravo2019use,lauri2016exploration,kim2021plgrim,vanegas2016uav,xiao2019objectsearch,pajarinen2022composition}
& $\times$
& △
& $\bigcirc$
& $\times$
& $\times$ \\


Symbolic planning under uncertainty

~\cite{li2026uncertainty,sanner2010symbolic,zhang2015corpp,serrano2021knowledge,
tang2025tru}
& $\bigcirc$
& $\bigcirc$
& $\bigcirc$
& $\times$
& $\times$ \\

Human-on-the-loop systems

~\cite{scharre2015introduction,
nahavandi2017trusted,
kelly2019hg,mandlekar2020human}
& $\times$
& $\times$
& △
& $\bigcirc$
& $\times$ \\

Uncertainty-aware query systems

~\cite{liang2024introspective,
mullen2025lbap}
& △
& $\times$
& $\bigcirc$
& $\bigcirc$
& $\bigcirc$ \\

KnowNo~\cite{ren2023robots}
& △
& $\times$
& $\bigcirc$
& $\bigcirc$
& $\bigcirc$ \\

\midrule
\textbf{Ours}
& $\bigcirc$
& $\bigcirc$
& $\bigcirc$
& $\bigcirc$
& $\bigcirc$ \\

\bottomrule
\end{tabularx}
\end{table*}



\subsection{Knowledge-based Robotic System}~\label{rel:kb}
In robotics, a knowledge base refers to a structured representation of information required for decision making. These systems use logic-based representations of objects, attributes, and relations to describe the world in a way that is understandable not only by robots but also by users~\cite{brachman2004knowledge, russell1995modern}. In particular, predicate-based symbolic representations such as~\cite{fikes1971strips, aeronautiques1998pddl} allow the system to describe both object properties and relationships among objects. In addition, graph-based models such as~\cite{johnson2015image, armeni20193d} are often used to represent these relationships more clearly by organizing objects and their connections in a simple structured form.


To systematically define robot states and actions, planning languages such as PDDL~\cite{aeronautiques1998pddl} or FDR~\cite{helmert2009concise} are widely adopted. In these representations, a state is defined as a set of predicates, and each action is described by its preconditions and effects, which specify how the state changes after execution. This formulation provides a clear model of state transitions and allows robots to plan over a well-defined symbolic state space. As a result, classical planning techniques~\cite{hoffmann2001ff, helmert2006fast, correa2024lifted} can be applied to compute action sequences that achieve a given goal.


Knowledge-based robotic systems store observed facts and infer new knowledge from existing information within this symbolic representation. Logical and relational inference allows robots to identify facts that are not directly observed, combine multiple relations, and obtain plausible hidden states from contextual knowledge~\cite{tenorth2012knowledge,liu2020graph}. For example, a robot can infer the likely location of an unseen object from qualitative spatial relations, container-location relations, or place-object co-occurrence knowledge~\cite{kunze2014using,hasegawa2023inferring}. Commonsense reasoning connects symbolic knowledge with perceptual evidence and supports decision making in partially observable environments~\cite{mota2021integrated,zhang2024icorpp,daruna2019robocse}. Ontology structures and type hierarchies further support consistent integration of new observations and generalization across related objects, actions, and task contexts~\cite{mitrevsk2021ontology,olivares2019review,manzoor2021ontology}.

While symbolic knowledge-based systems provide structured representations and support logical inference, recent approaches have explored the use of large language models (LLMs) to further expand how knowledge is utilized in robotic planning. Early LLM-based approaches employ procedural and commonsense knowledge encoded in LLMs to propose task decompositions or action sequences. The proposed plans are then grounded with robot affordances, programmatic interfaces, or execution feedback~\cite{ahn2022can,huang2022inner,singh2023progprompt,liang2023code,ding2023task}. Recent LLM-based approaches also incorporate external knowledge sources to compensate for incomplete observations. Some methods use scene graphs or open-vocabulary spatial representations to provide structured information about objects, places, and relations~\cite{gu2024conceptgraphs,leanza2025conceptbot}. Other methods retrieve demonstrations, procedural knowledge, or task-relevant constraints to guide long-horizon planning~\cite{singh2025adaptbot,zhouenhancing,xu2024p,wang2025instructrag,kumar2026open}. Related methods combine LLM-generated plans with symbolic constraints, spatial reasoning, or verification mechanisms to improve executability and consistency~\cite{huang2023voxposer,dai2024optimal,liu2025delta,paulius2025bootstrapping,quartey2025verifiably}.

However, many knowledge-based robotic systems still rely on deterministic symbolic representations and assume a single consistent world state, which limits their ability to represent multiple hypotheses under uncertainty~\cite{yang2018peorl,bouguerra2008monitoring}. As a result, such systems do not explicitly model uncertainty in action effects, sensing, or world states, and inferred facts are typically committed as definite knowledge. Recent LLM-based planning methods also suffer from grounding errors, incomplete retrieval, and hallucinated inferences, and do not maintain structured uncertainty over alternative states. While knowledge-based systems provide rich semantic structure for planning, they do not directly address uncertainty in observations, inferred facts, and knowledge validity. 


To address this limitation, we maintain multiple symbolic hypotheses over possible worlds instead of committing inferred facts as deterministic knowledge. Our framework combines symbolic knowledge-based reasoning with planning under uncertainty, allowing the robot to reason over uncertain inferred states while selectively requesting human feedback if the uncertainty becomes high.



\subsection{Planning under Uncertainty}~\label{rel:PlanningUncertainty}


\subsubsection{POMDP}~\label{rel:pu_pomdp}
Unlike deterministic symbolic knowledge-based systems that assume a single world state, POMDPs explicitly reason over uncertainty through belief representations over multiple possible states. In robotic environments, such uncertainty naturally arises from sensor noise, occlusions, incomplete observations, and unexpected state changes. Sequential decision-making under these conditions has been extensively studied within the framework of POMDPs~\cite{somani2013despot,bai2015intention,cramer2021probabilistic,bravo2019use}.

A key distinction between a POMDP and an MDP is the introduction of a belief state, which represents a probabilistic distribution over possible states when the true state cannot be directly observed. Accordingly, a wide range of solution methods have been proposed, including point-based value iteration~\cite{pineau2003point}, online tree search~\cite{ross2008online}, and Monte Carlo tree search (MCTS)~\cite{silver2010monte}.

Despite these advances, effectively representing and updating the belief state remains a fundamental challenge. Since a belief is inherently a continuous distribution over the state space, exact inference is computationally difficult. Moreover, as the state space grows, the belief space expands combinatorially, making inference increasingly intractable. As a result, various approaches have been proposed to approximate the belief state.

\subsubsection{Belief Approximation}~\label{rel:pu_belief_approx}
Studies on belief approximation can be broadly categorized into particle-based and factored-based. Particle-based methods approximate the belief as a set of sampled states and have been widely adopted due to their natural compatibility with sampling-based online planning~\cite{doshi2005particle, lim2023optimality}. For example,~\cite{silver2010monte} maintains a set of particles as a history-conditioned belief and estimates action values by sampling future trajectories from a generative model without performing exact Bayesian belief updates. Similar sampling-based ideas have also been extended to large and continuous POMDPs, where DESPOT~\cite{somani2013despot} evaluates a limited set of sampled scenarios and POMCPOW~\cite{sunberg2018pomcpow} controls search tree growth in continuous action and observation spaces.


In robotic domains, particle-based belief representations are widely used in localization, navigation, SLAM, object search, and manipulation, where particles encode hypotheses over robot, map, and object states under uncertainty. In navigation,~\cite{lauri2016exploration} combines POMCP with model predictive control to maximize information gain over occupancy-grid maps, while~\cite{kim2021plgrim} uses hierarchical local online POMDP planning for exploration in large unknown environments. Similar representations have also been used for continuous control under uncertainty, including autonomous driving at occluded intersections~\cite{brechtel2014continuous} and UAV collision avoidance or target tracking under partial observability~\cite{bai2011uav,vanegas2016uav}. In object search and manipulation, particles often represent complete scene hypotheses. For example,~\cite{xiao2019objectsearch} samples visible and occluded object configurations for cluttered object search, while~\cite{pajarinen2017multiobject,pajarinen2022composition} uses particle-based scene hypotheses for multi-object manipulation. Collectively, these studies show that particle beliefs provide a flexible approximation for preserving multimodal dependencies among robot, environment, and object states.



In contrast, factored belief representations model the belief by decomposing the state into a set of variables and exploiting their structure~\cite{hansen2000dynamic, feng2001approximate, veiga2014point}. Such approaches have been widely used in structured probabilistic models, including dynamic Bayesian networks and Hidden Markov Models~\cite{eddy1996hidden}, where the state is represented in a factored form. To improve computational efficiency, various approximate inference methods have been proposed to reduce the complexity of belief updates by relaxing dependencies among variables. For example,~\cite{murphy2013factored} approximates belief updates by maintaining a factorized representation while partially accounting for dependencies between variables. While factored representations exploit internal structure of state space, their independence assumptions break in symbolic planning where facts are tightly linked by action constraints. Therefore, we use a particle-based approximation to preserve these dependencies.


\subsubsection{Symbolic Planning under Uncertainty}~\label{rel:pu_sym}
Symbolic planning formalisms such as PDDL~\cite{aeronautiques1998pddl}, STRIPS~\cite{fikes1971strips}, and FDR~\cite{helmert2009concise} have been widely used due to their interpretability~\cite{ghallab2016automated}. These representations are typically based on first-order logic or similar discrete formalisms~\cite{lifschitz2019answer} which describes planning problems through explicit constraints and state transition rules. However, since traditional predicate logic only allows binary true or false values, it is difficult to directly model observation noise or state uncertainty.
To address this limitation, various approaches have attempted to integrate probabilistic reasoning into logical representations. For example, ~\cite{richardson2006markov} assign weights to logical rules to enable probabilistic inference over uncertain relationships, and~\cite{qu2019probabilistic} combine logical structure with probabilistic reasoning.

However, these approaches primarily focus on building inference models, and their high cost in model construction is unsuitable for online decision-making with complex constraints. 

As a result, several studies have explored planning under uncertainty while preserving symbolic structure~\cite{li2026uncertainty,sanner2010symbolic}. For example,~\cite{zhang2015corpp} encodes commonsense knowledge in ASP~\cite{lifschitz2019answer} to construct a prior over possible worlds and integrates it with probabilistic planning~\cite{kurniawati2008sarsop}. However, since the belief is defined over a set of possible worlds, the computational cost increases as the number of states grows. To address this, some works organize possible worlds in a structured manner. For instance,~\cite{serrano2021knowledge} construct a state space tree and a hierarchy of POMDPs based on domain knowledge, reducing complexity by leveraging abstraction. More recently, large language models have been used for belief approximation, as in~\cite{tang2025tru}, which builds a tree of hypotheses and uses LLMs to generate particles in open-ended settings. While these approaches enable more flexible or scalable belief representations, they either depend on predefined hierarchical structures or rely on the quality and consistency of generated hypotheses.

Prior work has explored various belief representations, and global belief approaches can be sample-intensive under strong symbolic constraints. Motivated by the broader idea of focusing computation on reachable belief regions~\cite{kurniawati2008sarsop}, we approximate belief locally by sampling states reachable under each action. These local beliefs are organized in a tree structure, and planning is performed using the off-the-shelf planner POMCP within an interpretable symbolic space. Real-world robotic systems also benefit from information that directly disambiguates competing hypotheses, which motivates human corrective systems that selectively request feedback under high uncertainty.


\subsection{Human Corrective System}~\label{rel:HCS}
\ny{Despite recent advances in robotic systems, robots still face operational uncertainty in dynamic and unpredictable real-world environments. Prior work has addressed this limitation by involving humans in different forms of robot autonomy~\cite{kelly2019hg, mandlekar2020human}. In particular, in human-on-the-loop (HOTL) systems, robots execute tasks autonomously, while human operators monitor the system and revise or override robot decisions if failures or errors occur~\cite{scharre2015introduction, nahavandi2017trusted}. This supervisory structure can improve safety and robustness because human operators can compensate for errors that autonomous systems cannot resolve.}

\ny{However, HOTL systems require operators to sustain attention and detect failures during robot operation. To provide effective corrective feedback, operators must first notice that the robot is in an uncertain or erroneous state. Situation awareness research suggests that operators must then understand the current situation, project possible consequences, and determine an appropriate intervention~\cite{endsley2018automation}. This cognitive process can be difficult in repetitive tasks. Operators may experience vigilance decrement because continuous exposure to monotonous monitoring tasks can slow responses and reduce detection accuracy~\cite{davies1982psychology, davies2013human}. Operators may also experience inattentional blindness, in which visible failures are missed if attention is not directed to the relevant event~\cite{mack2003inattentional}. As a result, continuous supervision can increase workload and fatigue~\cite{wong2017workload}. Continuous supervision can also reduce the likelihood that operators detect critical robot failures~\cite{wong2017workload, galster2006managing}.}

\ny{Prior work has shown that robots can improve task performance, team efficiency, and subjective satisfaction if they proactively provide assistance during human-robot collaboration~\cite{baraglia2017efficient}. This finding suggests that proactive robot behavior can support collaboration by reducing the need for humans to manage every decision or detect every need for intervention. In parallel, uncertainty-aware robot systems have advanced toward determining when human input is needed. For example, conformal prediction-based methods construct prediction sets with statistical guarantees and use the size of these sets to decide whether to query the user~\cite{ren2023robots, liang2024introspective, mullen2025lbap}. These methods show that robots can use uncertainty estimates to reduce unnecessary human intervention.}

\ny{However, studies that develop and evaluate system-initiated query systems have mainly focused on technical performance, such as task performance, uncertainty calibration, and query efficiency. Human factors, such as workload and fatigue, can also affect task performance in automated systems that involve human operators~\cite{murata2000ergonomics, borragan2017cognitive, rajavenkatanarayanan2019towards}. Therefore, system-initiated corrective interaction should be evaluated not only in terms of technical performance but also in terms of operator burden and fatigue.}

\ny{Building on this gap, we examine how system-initiated corrective interaction affects operator workload, fatigue, and task performance in automated robot systems. This motivates the user-study component of our evaluation and complements our technical contributions on belief-based feedback triggering and query selection. In our interaction structure, the robot monitors its own belief uncertainty and asks for corrective feedback when additional information is useful. Users therefore provide targeted input at system-identified moments rather than continuously monitoring for failures.}
